% F-from-v2.tex: что версия 3 изменила по сравнению с версией 2 и какие её утверждения отозвала.
\section{От версии 2 к версии 3}
\label{app:from-v2}

Версия 2 сохранена без изменений рядом с этим файлом. Настоящее приложение записывает, что между
ними изменилось, и, что полезнее, каких предложений версии 2 настоящая версия больше написать не
может.

\subsection*{Что приобрела система}

Между \code{в9.355} и \code{1.0.0-кв.1}, в том порядке, в каком их встречают кольца: основание,
на котором держится регистрация, доверенный поручитель и регистрационный код; карта как
кодировка, эмулятор, шаг персонализации и устройство проверки, с принуждением, измеренным во
времени; решения об органах, операторах, квотах, потолках и доверяющих сторонах, сохраняемые как
записи только на добавление; восемь мутационных учений над проверочным слоем и поправки, которых
они потребовали; четыре формальные модели, проверяемые при каждой отправке; работающая и
измеренная миграция населения; правила кэширования, позволяющие недоверенной сети переносить
состояние; конвейер эпох глубины 24; измеренное изложение того, что могут две сговорившиеся
доверяющие стороны; резервный регион; модель ёмкости, нашедшая целочисленное поле, и модель
стоимости; публичный отчёт о прозрачности; сворачивание пилота и план сосуществования;
сопоставление уровней достоверности, проверяемое машиной, и автоматизируемая треть доступности;
два формата-моста; четыре пакета в публичных реестрах; верификатор ОИД4ВП; операционный договор и
сводная таблица; и первые три внешних существительных в ней.

\subsection*{Утверждения версии 2, которые настоящая версия отзывает или сужает}

{\footnotesize
\setlength{\tabcolsep}{4pt}\renewcommand{\arraystretch}{1.12}
\rowcolors{2}{tabstripe}{white}
\begin{xltabular}{\textwidth}{>{\RaggedRight\arraybackslash\hspace{0pt}}p{2.0in} >{\RaggedRight\arraybackslash\hspace{0pt}}X >{\RaggedRight\arraybackslash\hspace{0pt}}p{1.6in}}
\caption{Предложения версии 2, которых версия 3 не повторяет.}\label{tab:withdrawn}\\
\toprule
\rowcolor{tabheader}\thead{Версия 2 говорила} & \thead{Версия 3 говорит} & \thead{Почему} \\
\midrule
\endfirsthead
\toprule
\rowcolor{tabheader}\thead{Версия 2 говорила} & \thead{Версия 3 говорит} & \thead{Почему} \\
\midrule
\endhead
\midrule\multicolumn{3}{r}{\itshape продолжение на следующей странице}\\
\endfoot
\bottomrule
\endlastfoot
Постквантовая, несвязываемая со стороны эмитента, устойчивая к принуждению система идентификационных токенов & Несвязываемая со стороны эмитента, осведомлённая о принуждении система идентификационных токенов, подписанная МЛ-ДСА-65 при проверяемом пути миграции алгоритмов & Два слова ослаблены собственными оценками репозитория; каждое теперь отвергается проверкой на внешних поверхностях \\
Алгоритм подписи по умолчанию постквантовый с первого дня & Алгоритм есть МЛ-ДСА-65, а путь разработки и НИ по умолчанию подписывает помеченной заглушкой; промышленный запуск без настоящего подписания падает закрытым & Измерено: 32 байта заглушки там, где подпись занимает 3309 \\
Путь принуждения есть проверенное свойство; внешняя сторона его побочный канал не измеряла & Механизм противостоит беспечному принуждающему, но не осведомлённому, стоящему рядом с держателем, и против законного доступа в итоге вредит & Оценка лаборатории против трёх принуждающих, названных договором \\
Настоящей реализации сторонней командой не случилось & Один неизменённый сторонний кошелёк предъявил по ОИД4ВП и был принят; реализации родного протокола кем-то ещё нет & Строка кошелька в сводной таблице; граница между двумя путями \\
Спецификация и набор проверок соответствия существуют; интеграция не доказана & Размещённый пакет Фонда ОпенАйДи прогнал профиль: 7 пройдено, 4 на рассмотрении; более поздний прогон против \code{1.0.0rc7} сертифицирован (24 сентября 2026) & Собственные подписанные выгрузки Фонда; его перечень сертифицированных реализаций \\
Межверификаторная соотносимость есть постоянное задокументированное свойство, затем ограниченное & Ограничена на пути с нулевым разглашением в пределах толпы эпохи против противника, сопоставляющего по равенству; открыта при обыкновенном предъявлении и предъявлении СД-ЖВТ; у толпы нет нижнего предела & Три находки лабораторной работы о связываемости \\
Сертификат, которым верификатор подписывает запросы, был правильным & У него не было назначения ключа, и каждый внутренний инструмент его пропустил & Его отверг сторонний кошелёк \\
Слой проверок ловит собственное нарушение & Семьдесят одна проверка из семидесяти пяти проходила на дереве, где их свойство было закомментировано, а тесты ограничений, триггеров, индексов и процедур часто не заметили бы исчезновения своего предмета & Мутационные учения; каждая находка закрыта и отработана \\
Тридцать шесть таблиц, 26 триггеров, 191 проверка, 48 случаев соответствия, 18 задач НИ, 20 форматов, 1\,048 функций тестов, 755 и 84 из 89 тестов & Сорок пять таблиц, 42 триггера, 272 проверки, 118 случаев, 20 задач, 27 форматов, 1\,755 функций тестов, 992 и 107 из 112 тестов & Перемерено на \code{1.0.0-кв.1}; приложение~\ref{app:counts} \\
Тринадцать случаев аудита как самой записи & Четырнадцать названных, тридцать одна таблица под триггером добавления & Собственные триггеры схемы, пересчитанные \\
Физический токен смоделирован, но не произведён & Закодирован, эмулирован, персонализирован и прочитан эталонным устройством; не произведён & \code{полярис\_карта/} и её учения \\
Восстановление между регионами не построено & Резервный регион построен, и его ненулевая точка восстановления измерена & Учение об эвакуации \\
Цели дорожной карты предстоит проверить & Проверены: каждая цель по пропускной способности перекрыта на порядок, а 32-битное поле остановило бы службу за пять дней & Модель ёмкости \\
Версия есть в9.355, счёт отправок & Версия есть \code{1.0.0-кв.7}, и сдвигается она лишь изменением, видимым снаружи & Правило версионирования операционного договора \\
Будущая работа есть перечень того, что предстоит построить & Будущая работа есть перечень того, что пришлось бы сделать внешним сторонам, вопросов, на которые лаборатория должна ответить, и пунктов, отложенных до тех пор, пока в них не возникнет нужда у кого-то извне & Правило слияния и рубеж прекращения операционного договора \\
\end{xltabular}
}

\subsection*{Что намеренно не сделано}

Стенд производительности и базовая линия приложения не перезапускались, и их строки в
таблице~\ref{tab:perf} несли свои прежние отметки; стенд перезапущен 23~сентября 2026 года
(приложение~\ref{app:counts}). Документы версий 1 и 2 не правились. Документы
репозитория, расходящиеся с кодом на \code{1.0.0-кв.7}, ради настоящего документа не
исправлялись; они перечислены в авторских заметках, как это было с версией 2, чтобы документ
следовал коду, а расхождение не потерялось.

\subsection*{Введённый словарь}

Пятая метка состояния, \sWit, для способности, которую задействовала названная внешняя сторона,
с правилом, что она ставится лишь там, где в сводной таблице есть заполненная строка, и никогда
не означает аудита, сертификации или пилотного внедрения. Тринадцатая ступень лестницы, для
знания, которого дерево не способно проверить о себе само. И собственные слова операционного
договора, внешнее существительное, сводная таблица, кандидат в выпуски, употребляемые так, как их
употребляет репозиторий.
